\styledchapter{Poisson Limit Theorem}

We eventually want to show that \[
	\operatorname{Bin}\left( n,p \right) \leadsto \operatorname{Poi}(\lambda)
\] when $np \to \lambda$.

\section{Coupling}

We introduce the notion of \emph{coupling}. Suppose 
\begin{align*}
	&X: (\Omega_1, \mathscr{F}_1, \P_1) \to \left( \R,\mathscr{B}\left( \R \right) \right) \\
	&Y: (\Omega_2, \mathscr{F}_2, \P_2) \to \left( \R, \mathscr{B}_\R \right)
.\end{align*}

It would be nice if $X$ and $Y$ lived on the same space.

\begin{definition}[Coupling]
	The pair $\left( \hat{X}, \hat{Y} \right)$ is a coupling of $(X,Y)$ if $\hat{X}$ and $\hat{Y}$ are on the same space and \[
		\hat{X} \overset{\text{d}}{=} X, \quad \hat{Y} \overset{\text{d}}{=}Y
	.\] 
\end{definition}

\begin{example}
	Suppose $X \sim N(0,1)$ and $Y \sim N(1,1)$. One coupling is \[
		\hat{X} = X, \quad \hat{Y} = X+1
	.\] 
\end{example}

There are many possible couplings for the above example. Are any couplings more notable than others?

\begin{definition}[Maximal coupling]
	We say $(\hat{X},\hat{Y})$ is a maximal coupling of $(X,Y)$ if \[
		\P\left[\hat{X} = \hat{Y}\right] 
	\] is maximized.
\end{definition}
Clearly, the previous example has $\P\left[\hat{X} = \hat{Y}\right] = 0 $. How do we construct the maximal coupling?

\begin{theorem}[Construction of maximal coupling, discrete RVs]
	Suppose $X,Y$ are both discrete and notate\boxednote{The same argument applied to continuous RVs if we work with densities; take 383.} \[
		\P\left[X = i\right] = p_i, \quad \P\left[Y = i\right]  = q_k
	.\] 
	The maxmimal coupling satisfies\boxednote{The LHS is called total variation affinity/distance.} \[
		\P\left[\hat{X} \ne \hat{Y}\right] = \sum\limits_{i}^{} p_i \wedge q_i
	.\] 
\end{theorem}

\begin{proof}
	Step I (Showing an upper bound).
	Define $\alpha = \sum_{i}^{} p_i \wedge q_i$. For any coupling $(\hat{X}, \hat{Y})$ of $(X,Y)$, define \[
		A = \left\{i: p_i < q_i\right\}
	.\] Then we can write 
	\begin{align*}
		\P\left[\hat{X} = \hat{Y}\right] 
		&= \P\left[\hat{X} = \hat{Y} \in A\right] + \P\left[\hat{X}= \hat{Y} \in A^{c}\right] \\
		&\le \P\left[\hat{X} \in A\right] + \P\left[\hat{Y} \in A^{c}\right] \\
		&= \P\left[X \in A\right] + \P\left[Y \in A^{c}\right] \\
		&= \sum\limits_{i\in A}^{} p_i + \sum\limits_{i\in A^{c}}^{} q_i \\
		&= \sum\limits_{i}^{} p_i \wedge q_i = \alpha
	.\end{align*}

	Step II (Showing the upper bound is achieved). 
	Define \[
		\begin{cases}
			b_i &= \frac{p_i \wedge q_i}{\alpha} \\
			c_i &= \frac{p_i - p_i \wedge q_i}{1-\alpha} \\
			d_i &= \frac{q_i - p_i \wedge q_i}{1-\alpha}
		\end{cases}
	.\] We see that all three of $b_i,c_i,d_i$ are probabilities because they are nonnegative and individually sum to $1$.\boxednote{The denominators are scaling factors.}
	
	Generate independent random variables $B,C,D$ such that \[
		\P\left[B = i\right] = b_i, \quad \P\left[C = i\right]  = c_i, \quad \P\left[D = i\right]  = d_i
	.\] Then generate a binary indicator $I \sim \operatorname{Bernoulli}(\alpha)$.

	If $I = 1$, then construct $\hat{X}= \hat{Y} = B$. If $I = 0$, then put $\hat{X} = C, \hat{Y} = D$.\boxednote{If $I=1$, which happens with probability $\alpha$, $\hat{X}$ and $\hat{Y}$ agree. Otherwise, with probability $1-\alpha$, they disagree, as if the probability of $C=i$, $c_i$ is nonzero, the probability od $D=i$, $d_i$ is zero.} We have constructed $\hat{X},\hat{Y}$ on the same space. We need to check that it is a coupling of $(X,Y)$ and that it is indeed maximal.\boxednote{When making a maximal coupling, we begin with marginal distributions of $X$ and $Y$ and want to find the joint distribution such that they are as dependent as possible.}

	We can write
	\begin{align*}
		\P\left[\hat{X} = i\right] 
		&= \P\left[\hat{X} =  i\mid I = 1\right] \P\left[I = 1\right]  + \P\left[\hat{X} = i \mid I = 0\right] \P\left[I = 0\right] \\
		&= \P\left[B=i\right] \P\left[I = 1\right] + \P\left[C= i\right] \P\left[I=0\right] \\
		&= \frac{p_i \wedge q_i}{\alpha}\alpha + \frac{p_i - p_i \wedge q_i}{1-\alpha}(1-\alpha) \\
		&= p_i = \P\left[X = i\right] 
	.\end{align*}
	Similarly, \[
		\P\left[\hat{Y} = i\right] = q_i = \P\left[Y = i\right] 
	.\] 

	Therefore, \[
		\P\left[\hat{X} = \hat{Y}\right] \ge \P\left[I = 1\right]  = \alpha
	,\] so $(\hat{X},\hat{Y})$ is maximal.
\end{proof}

\section{LeCam and Stein--Chen Poisson Limit Theorems}

\begin{definition}[Total variation]
	For random variables $P,Q$ on the same measure space $\left( \Omega, \mathscr{F} \right)$, define \[
		\operatorname{TV}(P,Q) = \sup_{B \in \mathscr{F}} \left| P(B) - Q(B) \right|
	.\] 
\end{definition}
\begin{theorem} \label{3-1-1}
	Suppose $P,Q$ are discrete random variables on $\left( \Omega, \mathscr{F} \right)$. The maximial coupling $(\hat{X},\hat{Y})$ of $(P,Q)$ satisfies \[
		\P\left[\hat{X} \ne \hat{Y}\right] = \operatorname{TV}(P,Q)
	.\] 
\end{theorem}

\begin{proof}
	It suffices to show $\operatorname{TV}(P,Q) = 1- \underbrace{\sum_{i}^{} p_i \wedge q_i}_{\mathclap{\text{total variation affinity}}}$.

	Again, let $A = \left\{i: p_i < q_i\right\}$. Since $A \in \mathscr{F}$, the definition of total variation gives \[
		\operatorname{TV}(P,Q) \ge \left|P(A) - Q(A)\right| = Q(A) - P(A)
	.\] For every $B$, 
	\begin{align*}
		\left| P(B) - Q(B) \right|
		&= \left|\sum\limits_{i\in B}^{} (p_i - q_i) \right| \\
		&= \left| \sum\limits_{i \in B \cap A}^{} \underbrace{(p_i - q_i)}_{\le 0} + \sum\limits_{i \in B\cap A^{c}}^{} \underbrace{(p_i - q_i)}_{\ge 0}\right| \\
		&= \left| \sum\limits_{i \in B \cap A^{c}}^{} (p_i - q_i) - \sum\limits_{i\in B \cap A}^{} (q_i - p_i) \right|\\
		&\le \max \left\{\sum\limits_{i \in B \cap A^{c}}^{} (p_i - q_i) , \sum\limits_{i \in B \cap A}^{} (q_i - p_i)\right\} \\
		&\le \max \left\{\sum\limits_{i\in A}^{} (p_i - q_i), \sum\limits_{i \in A}^{} q_i - p_i\right\} \\
		&= \max \left\{1-P(A) - (1- Q(A)), Q(A) - P(A)\right\} \\
		&= Q(A) - P(A)
	.\end{align*}
	Taking a sup over all $B \in \mathscr{F}$, we have \[
		\operatorname{TV}(P,Q) \le Q(A) - P(A)
	.\] Combining with the other direction we already showed, we have equality.

	Since 
	\begin{align*}
		\operatorname{TV}(P,Q) &= Q(A) - P(A) \\
		\sum\limits_{i}^{} p_i \wedge q_i &= \sum\limits_{i \in A}^{} p_i + \sum\limits_{i \in A^{c}}^{} q_i \\
										  &= P(A) + Q(A^{c})
	\end{align*}
	sum to $1$, we see that \[
		\operatorname{TV}(P,Q) + \sum\limits_{i}^{} p_i \wedge q_i = 1
	,\] and subtract $\sum_{i}^{} p_i \wedge q_i$ to conclude.
\end{proof}

\begin{theorem}[LeCam Poisson Limit Theorem]
	Suppose $X_i \sim \operatorname{Bernoulli}(p_i)$ independently for all $i = 1,\ldots,n$ and define $\lambda = \sum_{i=1}^{n} p_i$. Then\boxednote{Like other notions of distance, the total variation between two random variables is really the total variation between their distributions.} \[
		\operatorname{TV}\left( \sum\limits_{i=1}^{n} X_i, \operatorname{Poi}(\lambda) \right) \le \sum\limits_{i=1}^{n} p_i^2
	.\] 
\end{theorem}
\begin{example}
	Suppose $p_i = \frac{\lambda}{n}$. Then we have \[
		\operatorname{TV}\left( \operatorname{Bin}(n, \frac{\lambda}{n}), \operatorname{Poi}(\lambda) \right) \le \frac{\lambda^2}{n} \to 0
	.\] 
	The strength of this result is that it only relies on independence. We allow $p_i$ to vary, have an explicit bound on the total variation, and have that total variation goes to zero, which implies weak convergence.\footnote{Total variation is the distance in measures across all sets, including half intervals. So we have convergence in CDF and thus weak convergence.}
\end{example}
\begin{proof}
	Define $Z_i \sim \operatorname{Poi}(p_i)$ independently for all $i$.

	Then \[
		\sum\limits_{i=1}^{n} Z_i \sim \operatorname{Poi}(\sum\limits_{i=1}^{n} p_i ) = \operatorname{Poi}(\lambda)
	.\] So it suffices to bound the total variation between $\sum_{i}^{} X_i$ and $\sum_{i}^{} Z_i$. By Theorem (\ref{3-1-1}), this is equal to the probability that the maximal coupling variables are not equal.
	
	Let $\left( \hat{X}_i, \hat{Z}_i \right)$ be the maximal coupling of $\left( X_i,Z_i \right)$. Then by our theorem,
	\begin{align*}
		\P\left[\hat{X}_i = \hat{Z}_i\right] 
		&= 1- TV\left( X_i, Z_i \right)\\
		&= \sum\limits_{j}^{} \P\left[X_i = k\right]  \wedge \P\left[Z_i = k\right]  \\
		\P\left[Z_i = k\right] &= \frac{p_i^{k}}{k!} e^{-p_i}, \quad \P\left[X_i = k\right] = \begin{cases}
			p_i, \quad k= 1 \\ 1-p_i, \quad k= 0
		\end{cases} \\
		\implies \P\left[\hat{X}_i = \hat{Z}_i\right] 
							   &= \min \left\{1-p_i, e^{-p_i}\right\} + \min \left\{p_i, p_i e^{-p_i}\right\} \\
							   &= 1- p_i + p_i e^{-p_i}	\\
							   &\ge 1- p_i + p_i (1-p_i) \\
							   &= 1-p_i^2
	.\end{align*}
		Note that
		\[
			\left( \sum_{i=1}^{n} \hat{X}_i, \sum_{i=1}^{n} \hat{Z}_i \right)
		\]
		is a coupling of
		\[
			\left( \sum_{i=1}^{n} X_i, \sum_{i=1}^{n} Z_i \right).
		\]
		It need not be maximal, but it suffices for the first inequality in
		\begin{align*}
			\operatorname{TV}\left( \sum\limits_{i=1}^{n} X_i, \operatorname{Poi}(\lambda) \right)
			&= \operatorname{TV}\left( \sum\limits_{i=1}^{n} X_i, \sum\limits_{i=1}^{n} Z_i \right) \\
		&= \P\left[\widehat{\sum\limits_{i=1}^{n} X_i} \ne \widehat{\sum\limits_{i=1}^{n} Z_i}\right]  \\
		&\le \P\left[\sum\limits_{i=1}^{n} \hat{X}_i \ne \sum\limits_{i=1}^{n} \hat{Z}_i\right] \\
		&\le \sum\limits_{i=1}^{n} \P\left[\hat{X}_i \ne \hat{Z}_i\right] \\
		&\le \sum\limits_{i=1}^{n} p_i^2
	.\end{align*}
\end{proof}

\begin{lemma}[Stein-Chen]
	We\boxednote{The forward direction lets $f$ be any function such that the expectations exist. A minimal assumption $\left| Zf(Z) \right|< \infty$. The backwards direction requires the identity hold for all bounded $f: \N_0 \to \R$.} have that \[
		Z \sim \operatorname{Poi}(\lambda) \iff \forall\ f, \E\left[Zf(Z)\right] = \lambda \E\left[f(Z+1)\right] 
	.\] 
\end{lemma}
\begin{proof}
	($\implies$) Since we have for every $k = 0,1,2,\ldots$ \[
		\P\left[Z = k\right] = \frac{\lambda^k}{k!} e^{-\lambda}
	,\] we know
	\begin{align*}
		\E\left[Zf(Z)\right]
		&= \sum\limits_{k=0}^{\infty} kf(k) \frac{\lambda^k}{k!} e^{-\lambda} \\
		&= \sum\limits_{k=1}^{\infty} k f(k) \frac{\lambda^k}{k!} e^{-\lambda} \\
		&= \lambda \sum\limits_{k=1}^{\infty} f(k) \frac{\lambda^{k-1}}{(k-1)!} e^{-\lambda} \\
		&= \lambda \sum\limits_{k=0}^{\infty} f(k+1) \frac{\lambda^k}{k!} e^{-\lambda} = \lambda \E\left[f(Z+1)\right] 
	.\end{align*}

	($\impliedby$) We use a difference equation instead of the differential equation in Stein's method. Consider for each $k=0,1,2,\ldots$ and the equation 
	\begin{equation} \label{2-26-1}
		\lambda f(k+1) - kf\left( k \right) = \mathds{1}_{k \in A} - \P\left[Z \in A\right] 
	\end{equation}for some set $A$.
	If\boxednote{If $W$ satisfies the Stein--Chen identity for every $f$, then plugging in the solution $f_A$ forces $\P(W\in A)=\P(Z\in A)$ for every set $A$, hence $W\sim\operatorname{Poi}(\lambda)$.} we can show that this equation has a solution $f_A$, then we know for a random variable $W$ with distribution satisfying the Stein--Chen identity for all $f$, we can plug in $f = f_A$ to get \[
		\P\left[W \in A\right] = \P\left[Z \in A\right] 
	,\] so $W \sim \operatorname{Poi}(\lambda)$.

	We would like to get equation (\ref{2-26-1}) into a form of $g(k+1) - g(k) = [\cdots]$, since then we can sum the right hand side to get $g$. Multiplying each side of (\ref{2-26-1}) by $\frac{\lambda^k}{k!}$, we have that for $k=1,2,\ldots$,
	\begin{equation}
		\frac{\lambda^{k+1}}{k!} f(k+1) - \frac{\lambda^k}{(k-1)!} f(k) = \frac{\lambda^k}{k!} \left( \mathds{1}_{k \in A} - \P\left[Z \in A\right]  \right)
	,\end{equation}
	and for $k=0$, \[
		\lambda f(1) = \mathds{1}_{0 \in A} - \P\left[Z \in A\right] 
	.\] 

	For $g(k) = \frac{\lambda^k}{(k-1)!} f(k)$, we have that \[
		\begin{cases}
			g(k+1) - g(k) = \frac{\lambda^k}{k!} \left( \mathds{1}_{ k \in A} - \P\left[Z \in A\right]  \right), \quad k= 1,2,\ldots \\
			g(1) = \mathds{1}_{0 \in A} - \P\left[Z \in A\right] 
		\end{cases}
	,\] 
	so we sum to get \[
		g(k+1) = \sum\limits_{j=0}^{k} \frac{\lambda^j}{j!} \left( \mathds{1}_{j \in A} - \P\left[Z \in A\right]  \right), \quad k = 0,1,2,\ldots
	.\] Therefore, our solution to the original difference equation is \[
		f_A(k+1) = \frac{k!}{\lambda^{k+1}} \sum\limits_{j=0}^{k} \frac{\lambda^j}{j!} \left( \mathds{1}_{j \in A} - \P\left[Z \in A\right]  \right), \quad k = 0, 1,2,\ldots
	.\] 
	Note that we don't actually need the solution $f(0)$, since the right hand side of the goal equation does not involve $f\left(0 \right)$ and the left hand side multiplies it by $0$.\boxednote{Review why the term $\mathds{1}_{j \in A} - \P\left[Z \in A\right] $ is necessary (the sum must be $0$).}

	Let $Z \sim \operatorname{Poi}(\lambda)$. Suppose that for some random variable $W$, for any $f$, \[
		\E\left[W f(W)\right] = \lambda \E\left[f(W+1)\right] 
	.\] 
	Plugging in our $f_A$, we know
	\begin{align*}
		0
		&= \E\left[\lambda f_A(W+1) - W f_A(W)\right] \\
		&= \E\left[\mathds{1}_{W \in A} - \P\left[Z \in A\right] \right] \\
		&= \P\left[W \in A\right] - \P\left[Z \in A\right]
	.\end{align*}
	Since $A$ is arbitrary, we know $W \sim \operatorname{Poi}(\lambda)$.
\end{proof}

\begin{lemma}
	The solution $f_A$ to the Stein-Chen equation is
	\[
		f_A(k+1) = \frac{k!}{\lambda^{k+1}} \sum\limits_{j=0}^{k} \frac{\lambda^j}{j!} \left( \mathds{1}_{j \in A}- \P\left[Z \in A\right]  \right)
	,\]
	and it satisfies, for integers $k,l$,
	\[
		\left| f_A(k) - f_A(l) \right| \le \min \left\{1,\frac{1}{\lambda}\right\} \left| k - l \right|
	.\]
\end{lemma}
\begin{proof}
	The formula for $f_A$ was established in the proof of the $(\impliedby)$ direction of the Stein--Chen lemma above. It remains to prove the Lipschitz bound.

	Write $f_A(k+1) = \sum_{i \in A} f_i(k+1)$, where $f_i := f_{\{i\}}$. Substituting $A = \{i\}$ into the formula and using the Poisson probability $\P[Z = j] = \frac{\lambda^j}{j!}e^{-\lambda}$,
	\begin{align*}
		f_i(k+1)
		&= \frac{k!}{\lambda^{k+1}} \sum\limits_{j=0}^{k} \frac{\lambda^j}{j!}\bigl(\mathds{1}_{j=i} - \P[Z=i]\bigr) \\
		&= \frac{k!}{\lambda^{k+1}} \cdot \frac{\lambda^i}{i!}\bigl(\mathds{1}_{i \le k} - \P[Z \le k]\bigr) \\
		&= \frac{\P[Z=i]}{\lambda \P[Z=k]}\bigl(\mathds{1}_{i \le k} - \P[Z \le k]\bigr),
	\end{align*}
	where the second line uses $\sum_{j=0}^{k}\frac{\lambda^j}{j!}\mathds{1}_{j=i} = \frac{\lambda^i}{i!}\mathds{1}_{i \le k}$ and \[\P[Z=i]\sum_{j=0}^{k}\frac{\lambda^j}{j!} = \frac{\lambda^i}{i!}e^{-\lambda}\cdot e^{\lambda}\P[Z\le k] = \frac{\lambda^i}{i!}\P[Z\le k],\] and the third uses \[
		\frac{k!}{\lambda^{k+1}}\cdot\frac{\lambda^i}{i!} = \frac{e^{-\lambda}}{\lambda \P[Z=k]}\cdot\frac{\P[Z=i]}{e^{-\lambda}} = \frac{\P[Z=i]}{\lambda \P[Z=k]}
	.\]

	We claim:
	\begin{enumerate}[label=(\roman*)]
		\item $f_k(k+1) \ge f_k(k)$,
		\item For $i \ne k$, $f_i(k+1) \le f_i(k)$.
	\end{enumerate}
	With these, we could bound $f_A(k+1) - f_A(k)$ from above: adding $k$ to $A$ increases the sum by (i), and each remaining $i \ne k$ term is non-positive by (ii), so
	\begin{align*}
		f_A(k+1) - f_A(k)
		&\le f_{A \cup \{k\}}(k+1) - f_{A \cup \{k\}}(k) \\
		&\le f_k(k+1) - f_k(k).
	\end{align*}
	We evaluate using $\frac{\P[Z=k]}{\P[Z=k-1]} = \frac{\lambda}{k}$:
	\begin{align*}
		f_k(k+1) &= \frac{\P[Z > k]}{\lambda}, \\
		f_k(k)   &= -\frac{\P[Z=k] \P[Z \le k-1]}{\lambda \P[Z=k-1]} = -\frac{\P[Z \le k-1]}{k},
	\end{align*}
	so
	\[
		f_k(k+1) - f_k(k) = \frac{\P[Z > k]}{\lambda} + \frac{\P[Z \le k-1]}{k}.
	\]
	For each $0 \le j \le k-1$, the Poisson recurrence gives $\P[Z=j] = \frac{j+1}{\lambda}\P[Z=j+1]$, so $\frac{\P[Z=j]}{k} = \frac{j+1}{k}\cdot\frac{\P[Z=j+1]}{\lambda} \le \frac{\P[Z=j+1]}{\lambda}$ since $j+1 \le k$. Summing,
	\[
		\frac{\P[Z \le k-1]}{k} \le \frac{1}{\lambda}\sum\limits_{j=0}^{k-1}\P[Z=j+1] = \frac{\P[1 \le Z \le k]}{\lambda}.
	\]
	Therefore,
	\[
		f_k(k+1) - f_k(k)
		\le \frac{\P[Z > k] + \P[1 \le Z \le k]}{\lambda}
		= \frac{\P[Z \ge 1]}{\lambda}
		= \frac{1 - e^{-\lambda}}{\lambda}.
	\]
	Since $1-e^{-\lambda} \le 1$ and $e^{-\lambda} \ge 1-\lambda$ gives $1-e^{-\lambda} \le \lambda$,
	\[
		f_k(k+1) - f_k(k) \le \frac{1-e^{-\lambda}}{\lambda} \le \min\!\left\{1,\frac{1}{\lambda}\right\}.
	\]

	For the other direction, note that $\P[Z \in \mathbb{N}_0] = 1$, so applying the formula gives $f_{\mathbb{N}_0} \equiv 0$ and thus $f_{A^c} = -f_A$. Thus
	\[
		f_A(k) - f_A(k+1) = f_{A^c}(k+1) - f_{A^c}(k) \le \min\!\left\{1,\frac{1}{\lambda}\right\},
	\]
	by the same argument applied to $A^c$. Combining both directions,
	\[
		\bigl|f_A(k+1) - f_A(k)\bigr| \le \min\!\left\{1,\frac{1}{\lambda}\right\},
	\]
	and the triangle inequality extends this to $|f_A(k) - f_A(l)| \le \min\{1,\tfrac{1}{\lambda}\}|k - l|$.

	It remains to prove (i) and (ii). From the formulas above, $f_k(k+1) = \frac{\P[Z>k]}{\lambda} \ge 0$ and $f_k(k) = -\frac{\P[Z\le k-1]}{k} \le 0$, so $f_k(k+1) \ge 0 \ge f_k(k)$, proving (i).

	For (ii), when $i < k$ both numerators involve the tail:
	\begin{align*}
		f_i(k+1) &= \frac{\P[Z=i] \P[Z > k]}{\lambda \P[Z=k]}, &
		f_i(k)   &= \frac{\P[Z=i] \P[Z > k-1]}{\lambda \P[Z=k-1]},
	\end{align*}
	so $f_i(k+1) \le f_i(k)$ iff $\frac{\P[Z>k]}{\P[Z=k]} \le \frac{\P[Z>k-1]}{\P[Z=k-1]}$. Writing
	\[
		\frac{\P[Z > k]}{\P[Z=k]} = \sum\limits_{l > 0} \lambda^l \frac{k!}{(l+k)!},
	\]
	each term $\frac{k!}{(l+k)!} = \frac{1}{(k+1)\cdots(k+l)}$ is decreasing in $k$, so the sum is decreasing in $k$, which gives $f_i(k+1) \le f_i(k)$.

	When $i > k$ both numerators involve the CDF:
	\begin{align*}
		f_i(k+1) &= -\frac{\P[Z=i] \P[Z \le k]}{\lambda \P[Z=k]}, &
		f_i(k)   &= -\frac{\P[Z=i] \P[Z \le k-1]}{\lambda \P[Z=k-1]},
	\end{align*}
	so $f_i(k+1) \le f_i(k)$ iff $\frac{\P[Z\le k]}{\P[Z=k]} \ge \frac{\P[Z\le k-1]}{\P[Z=k-1]}$. Writing
	\[
		\frac{\P[Z \le k]}{\P[Z=k]} = \sum\limits_{l=0}^{k} \lambda^{-l} \frac{k!}{(k-l)!},
	\]
	both the number of terms and each factor $\frac{k!}{(k-l)!} = k(k-1)\cdots(k-l+1)$ increase in $k$, so the ratio is increasing in $k$, giving $f_i(k+1) \le f_i(k)$.
\end{proof}

\begin{proposition}
	For $Z \sim \operatorname{Poi}(\lambda)$ and a random variable $W$, the total variation is
	\begin{align*}
		\TV\left( W,Z \right)
		&= \sup_A \left| \P\left[W \in A\right]  - \P\left[Z \in A\right]  \right| \\
		&= \sup_A \left| \E\left[W f_A(W) \right] - \lambda \E\left[f_A(W+1)\right]  \right| \\
		&\le \sup_{\operatorname{Lip}(f) \le \min\left\{1, \frac{1}{\lambda}\right\}} \left| \E\left[W f(W) - \lambda\E\left[f(W+1)\right] \right]  \right|
	.\end{align*}
\end{proposition}

\begin{proof}
	The first equality is the definition of total variation, the second comes from the Stein--Chen identity, and the third from the Lipschitz bound we just derived.
\end{proof}

\begin{theorem}[Stein-Chen]
	Suppose we have independent $X_i \sim \operatorname{Bernoulli}(p_i)$ for $i=1,\ldots,n$. Put $\lambda = \sum_{i=1}^{n} p_i$. Then we have \[
		\TV\left( \sum_{i=1}^{n} X_i, \operatorname{Poi}(\lambda) \right) \le \min \left\{1,\frac{1}{\lambda}\right\} \sum_{i=1}^{n} p_i^2
	.\] 
\end{theorem}

We have a stronger result than LeCam's result because of the min term. When $\lambda\to \infty$ and $\sum_{i=1}^{n} p_i^2$ is controlled, we get that $\sum_{i=1}^{n}X_i \leadsto \operatorname{Poi}(\lambda)$.
\begin{proof}
	Let $W = X_1 + \cdots + X_n$, with $X_i \sim \operatorname{Bernoulli}(p_i)$ independent for $i = 1, \ldots, n$. Write $W_i = W - X_i$ and put $\lambda = \sum_{i=1}^{n} p_i$.

	Step I (Decomposing the discrepancy). We decompose the Stein--Chen discrepancy as
	\begin{align*}
		\lambda \E\left[f(W+1)\right] - \E\left[Wf(W)\right]
		&= \sum\limits_{i=1}^{n} p_i \E\left[f(W+1)\right] - \sum\limits_{i=1}^{n} \E\left[X_i f(W_i + X_i)\right]
	.\end{align*}
	Since $X_i \in \{0,1\}$ and $X_i \ind W_i$, conditioning on $W_i$ yields
	\begin{align*}
		\E\left[X_i f(W_i + X_i)\right]
		&= \E\left[\E\left[X_i f(W_i + X_i) \mid W_i\right]\right] \\
		&= \E\left[p_i \cdot f(W_i + 1) + (1-p_i) \cdot 0\right] \\
		&= p_i \E\left[f(W_i + 1)\right]
	,\end{align*}
	where the $X_i = 0$ term vanishes because $X_i \cdot f(W_i + X_i) = 0 \cdot f(W_i) = 0$ when $X_i = 0$. Substituting back,
	\begin{align*}
		\lambda \E\left[f(W+1)\right] - \E\left[Wf(W)\right]
		&= \sum\limits_{i=1}^{n} p_i \left(\E\left[f(W+1)\right] - \E\left[f(W_i+1)\right]\right)
	.\end{align*}

	Step II (Applying the Lipschitz bound). Since $W = W_i + X_i$ and $X_i \in \{0,1\}$, we have $|W - W_i| = X_i$. The Lipschitz bound on $f_A$ from the lemma above then gives $\left|f(W+1) - f(W_i+1)\right| \le \min\left\{1,\frac{1}{\lambda}\right\} X_i$, so
	\begin{align*}
		\left|\lambda \E\left[f(W+1)\right] - \E\left[Wf(W)\right]\right|
		&\le \sum\limits_{i=1}^{n} p_i \E\left[\left|f(W+1) - f(W_i+1)\right|\right] \\
		&\le \min\left\{1,\frac{1}{\lambda}\right\} \sum\limits_{i=1}^{n} p_i \E\left[X_i\right] \\
		&= \min\left\{1,\frac{1}{\lambda}\right\} \sum\limits_{i=1}^{n} p_i^2
	.\end{align*}
	The total variation bound then follows from the proposition above.
\end{proof}

\begin{example}
	Consider the typical example $p_i = \frac{\lambda}{n}$. Then we get the Law of Small Numbers.
\end{example}

\begin{proposition}[TV convergence implies weak convergence]
	If $\mu_n \to \mu$ in total variation, i.e.\ $\TV(\mu_n, \mu) \to 0$, then for every bounded measurable function $f$,
	\[
		\int f\,d\mu_n \to \int f\,d\mu.
	\]
	In particular, $\mu_n \leadsto \mu$.
\end{proposition}

\begin{proof}
	For any set $A \in \mathscr{F}$,
	\[
		\left|\int \mathbf{1}_A\,d\mu_n - \int \mathbf{1}_A\,d\mu\right| = |\mu_n(A) - \mu(A)| \le \TV(\mu_n,\mu) \to 0,
	\]
	so the result holds for indicators. By linearity it extends to simple functions. For general bounded measurable $f$, given $\varepsilon > 0$ choose a simple function $s$ with $\|f - s\|_\infty < \varepsilon$. Then
	\begin{align*}
		\left|\int f\,d\mu_n - \int f\,d\mu\right|
		&\le \left|\int (f-s)\,d\mu_n\right| + \left|\int s\,d\mu_n - \int s\,d\mu\right| + \left|\int (s-f)\,d\mu\right| \\
		&\le 2\varepsilon + \left|\int s\,d\mu_n - \int s\,d\mu\right|.
	\end{align*}
	The second term tends to $0$ as $n \to \infty$. Since $\varepsilon$ was arbitrary, $\int f\,d\mu_n \to \int f\,d\mu$. Weak convergence follows because bounded continuous functions are bounded and measurable.
\end{proof}

\begin{example}
	The boundedness assumption cannot be dropped. Let $\mu_n$ assign mass $\tfrac{1}{n}$ to $\{n\}$ and $1 - \tfrac{1}{n}$ to $\{0\}$, and let $\mu(x) = \mathds{1}_{0}$. The only points with positive probability are $0$ and $n$, so $\TV(\mu_n, \mu) = \frac{1}{n} \to 0$. However, taking $f(x) = x$ shows
	\[
		\int f\,d\mu_n = n \cdot \frac{1}{n} + 0 \cdot \left(1 - \frac{1}{n}\right) = 1 \ne 0 = \int f\,d\mu
	\]
	for every $n$, so $\int f\,d\mu_n \not\to \int f\,d\mu$.
\end{example}

\section{Dependence}


What if $X_1,\ldots,X_n$ are dependent?

\begin{definition}[Dependency neighborhood]
	For each $i \in \N$, consider $\left\{N_i\right\}$ where \[
		i \in N_i \subseteq [n] := \left\{1,2,\ldots,n\right\}
	\] such that \[
		X_i \ind \left\{X_j\right\}_{j \not\in N_i}
	.\] We call $N_i$ the dependency neighborhood of $X_i$.
\end{definition}

\begin{theorem}[Stein-Chen, dependent variables] \label{3-3-1}
	Suppose $X_i \sim \operatorname{Bernoulli}(p)$ for $i=1,\ldots,n$, where the $X_i$s are not necessarily independent. Put $\lambda = \sum_{i=1}^{n} p_i$. Then for dependency neighborhoods $N_i$, we have the total variation bound \[
		\TV\left( \sum\limits_{i=1}^{n} X_i, \operatorname{Poi}(\lambda) \right)
		\le \min\left\{1,\frac{1}{\lambda}\right\} \left( \sum\limits_{i=1}^{n} \sum\limits_{j \in N_i \setminus \left\{i\right\}}^{} \E\left[X_i X_j\right] + \sum\limits_{i=1}^{n} \sum_{j \in N_i}p_ip_j \right)
	.\] 
\end{theorem}
\begin{proof} 
	Fix a Stein solution $f$ with
	\[
		\left| f(k+1) - f(k) \right| \le \min\left\{1,\frac{1}{\lambda}\right\}.
	\]
	For each $i$, write
	\begin{align*}
		W
		&= X_i + \sum\limits_{j\in N_i \setminus \left\{i\right\}}^{} X_j + \sum\limits_{j \not\in N_i}^{} X_j \\
		&=: X_i + \sum\limits_{j \in N_i \setminus \left\{i\right\}}^{} X_j + V_i.
	\end{align*} 
	By assumption, $X_i \ind V_i$.

	Now decompose
	\begin{align*}
		\E\left[W f(W)\right] - \lambda\E\left[f(W+1)\right] 
		&= \sum\limits_{i=1}^{n} \left( \E\left[X_i f(W)\right] - p_i \E\left[f(W+1)\right]  \right)\\
		&= \text{(I)} + \text{(II)},
	\end{align*}
	where
	\begin{align*}
		\text{(I)}
		&= \sum\limits_{i=1}^{n} \left( \E\left[X_i f(W)\right] - \E\left[X_i f(V_i + 1)\right]  \right), \\
		\text{(II)}
		&= \sum\limits_{i=1}^{n} \left( \E\left[X_i f(V_i + 1)\right] - p_i \E\left[f(W+1)\right] \right).
	\end{align*}

	Since $X_i \in \{0,1\}$,
	\[
		\E\left[X_i f(V_i + X_i)\right]
		= \E\left[X_i f(V_i+1)\right]
		= p_i \E\left[f(V_i + 1)\right]
	,\] simplifying (II).

	Writing $f(W) - f(V_i +1) \le \left| f(W) - f(V_i+ 1) \right|$,
	\begin{align*}
		\left| \text{(I)} \right|
		&\le \sum\limits_{i=1}^{n} \E\left[X_i \left| f(W) - f(V_i + 1) \right|\right] \\
		&\le \min\left\{1,\frac{1}{\lambda}\right\}
			\sum\limits_{i=1}^{n}
			\E\left[X_i \sum\limits_{j \in N_i \setminus \left\{i\right\}}^{} X_j\right] \\
		&= \min \left\{1,\frac{1}{\lambda}\right\}
			\sum\limits_{i=1}^{n} \sum\limits_{j \in N_i \setminus \left\{i\right\}}^{} \E\left[X_i X_j\right]
	,\end{align*}
	where the second inequality follows from the Lipschitz bound on $f$.
	Also,
	\begin{align*}
		\left| \text{(II)} \right|
		&\le \sum\limits_{i=1}^{n} p_i \left| \E\left[f(V_i+1)\right] - \E\left[f(W+1)\right] \right| \\
		&\le \min \left\{1,\frac{1}{\lambda}\right\}
			\sum\limits_{i=1}^{n} p_i \E\left[\left| W - V_i \right|\right] \\
		&= \min \left\{1,\frac{1}{\lambda}\right\}
			\sum\limits_{i=1}^{n} p_i \E\left[\sum\limits_{j \in N_i}^{} X_j\right] \\
		&= \min \left\{1,\frac{1}{\lambda}\right\}
			\sum\limits_{i=1}^{n} \sum\limits_{j \in N_i}^{} p_i p_j.
	\end{align*}
	This gives the stated bound.
\end{proof}

\begin{example}
	Consider the undirected Erd\"os--R\'enyi graph with $n$ nodes and edges $A_{ij} \in \left\{0,1\right\}$ with $A_{ij} = A_{ji} \overset{\text{iid}}{\sim} \operatorname{Bernoulli}(p)$.

	We can count the number of edges with the CLT (for large $n$) or the independent version of the PLT (small $n$).

	More interestingly, we can count the number of triangles: \[
		T := \sum\limits_{1 \le i < j < k \le n}^{} A_{i j}A_{jk}A_{ki}
	.\] 
		We can use the dependency neighborhood condition. To make indexing easier, consider the dependency neighborhoods $N_\alpha$, so $\alpha \in N_\alpha$ and $X_\alpha \ind \left\{X_\beta\right\}_{\beta \not\in  N_\alpha}$. The dependency version of Stein-Chen PLT then says for $X_\alpha \sim \operatorname{Bernoulli}(p_\alpha)$ and $\lambda = \sum_\alpha p_\alpha$, we have \[
			\begin{aligned}
				\TV\left( \sum\limits_{\alpha}^{} X_\alpha, \operatorname{Poi}(\lambda) \right)
				&\le \min \left\{1,\frac{1}{\lambda}\right\} \sum\limits_{\alpha}^{} \sum\limits_{\beta \in N_\alpha \setminus \left\{\alpha\right\}}^{} \E\left[X_\alpha X_\beta\right] \\
				&\quad + \min \left\{1,\frac{1}{\lambda}\right\} \sum\limits_{\alpha}^{} \sum\limits_{\beta \in N_\alpha}^{} p_\alpha p_\beta
			\end{aligned}
		.\] 
	Write $\alpha = \left( i,j,k \right)$, so $X_\alpha = A_{ij}A_{jk}A_{ki}$. Since a triangle with a shared node shares exactly one edge, and each of the three edges of a triangle can be shared with triangles constructed by connecting the two nodes of the edge with $n-3$ different nodes, we see that when $N_\alpha$ are as small as possible, $\left| N_\alpha \right| = 3(n-3) +1$. (The dependency neighborhood of a triangle includes itself.)

	We see
	\begin{align*}
		\text{(II)}
		&= \sum\limits_{\alpha}^{} \sum\limits_{\beta \in N_\alpha}^{} p_\alpha p_\beta \\
		&= O\left( n^3 \cdot n \cdot p^6 \right) \\
		&= O\left( n^4 p^6 \right) \\
		\text{(I)}
		&= \sum\limits_{\alpha}^{} \sum\limits_{\beta \in N_\alpha\setminus \left\{\alpha\right\}}^{} \E\left[X_\alpha X_\beta\right]  \\
		&= \sum\limits_{\alpha}^{} \sum\limits_{\beta \in N_\alpha \setminus \left\{\alpha\right\}}^{} \E\left[A_{ab}A_{bc}A_{cd}A_{de}A_{ea}\right] \\
		&= O\left( n^4 p^5 \right) \\
		\lambda &= \binom{n}{3}p^3 = O\left( n^4 p^5, np^2\right)
	.\end{align*}
	Therefore, \[
		\TV\left( T, \operatorname{Poi}\left( \binom{n}{3}p^3 \right) \right) = O\left( \min \left\{n \cdot n^3 p^3, np^2\right\} \right)
	,\] 
	which tends to $0$ when $p \ll \frac{1}{\sqrt{n} }$.
\end{example}

\begin{remark}
	We can prove this dependent PLT using the martingale CLT instead of Stein's method.
\end{remark}



